\section{Magisterial chronology}
\label{app:magisterium}

The acts below establish different kinds of claims. The medieval councils'
confession of creation, modern doctrinal exposition, and the regulation of
devotion have their own objects and authority.

\begin{longtable}{@{}>{\raggedright\arraybackslash}p{0.14\linewidth}>{\raggedright\arraybackslash}p{0.30\linewidth}>{\raggedright\arraybackslash}p{0.49\linewidth}@{}}
\toprule
Date & Act and locus & Teaching relevant to the angels\\
\midrule\endhead
1215 & Lateran IV, \emph{Firmiter credimus}, DS~800--801 &
One Creator of spiritual and corporeal creatures; demons created good,
becoming evil through their own action; perpetual punishment with the devil.
\textbf{Defined}.\\
24 April 1870 & Vatican I, \emph{Dei Filius}, ch.~1 and Creator canons~1--5 &
Free creation from nothing; God's substance distinct from created spiritual
and bodily things; all creation governed by providence.
\textbf{Defined}.\\
12 August 1950 & Pius XII, \emph{Humani generis}, no.~26 &
Rejection of doubts about angelic personality and the essential distinction
of matter and spirit; defense of the gratuity of supernatural calling.
\textbf{Church teaching}.\\
1992; Latin typical edition 1997 & \emph{Catechism}, 328--336, 391--395 &
Personal spiritual creatures, their service of Christ and the Church,
the free angelic fall and the limits of Satan's power.
\textbf{Church teaching}.\\
17 December 2001 decree & Congregation for Divine Worship,
\emph{Directory on Popular Piety and the Liturgy}, 213--217 &
Angelic devotion ordered to the Church's faith and worship; restriction of
the discouragement of assigning other names than Michael, Gabriel and Raphael.
\textbf{Church teaching and devotional discipline}.\\
\bottomrule
\end{longtable}

The dating of the \emph{Catechism} identifies the promulgated work and its
typical edition; the English wording consulted is the Holy See's web witness
accessed on 29 September 2026. The Directory's final entry concerns
ecclesiastical direction of devotion; it does not constitute a new definition
of angelic nature.
